% chapters/chapter1/sec1.tex

\section{The essence, evolution and core characteristics of ISAC}
\label{sec:ch1-essential-definition}

% TODO: 本节由 Huixin Man 负责撰写。
% 写作目标：
% 1. 介绍无线通信系统的发展背景；
% 2. 引出 6G 对通信、感知、智能化融合的需求；
% 3. 说明 ISAC 的研究意义。

\subsection{Why is it necessary to redefine ISAC}
\label{subsec:ch1-why-redefine-isac}
\subsubsection{The long-standing separation of traditional "communication systems" and "radar systems"}
\label{subsubsec:ch1-traditional-separation}
无线通信系统与雷达系统是现代电子信息技术领域中最重要的两大分支。从工程历史上看，通信系统与雷达系统虽然都属于电磁信息系统，都依赖无线信号的发射、传播与接收，但它们自诞生之初就服务于不同任务，因此形成了两套相对独立的技术逻辑。自19世纪末无线电通信技术兴起以来，无线通信系统的核心目标始终未变——在发射端与接收端之间实现可靠、高效的信息传输。从经典信息论的视角来看，在加性高斯白噪声（AWGN）信道模型下，Shannon 信息论揭示了信道容量与系统带宽和信噪比之间的基本关系\cite{Shannon1948}:
\begin{equation}
C = B \log_2 \left( 1 + \frac{S}{N} \right)
\end{equation}
其中 $C$ 的单位是(bit/s)，$B$ 为带宽（Hz），$S$ 和 $N$ 分别为信号与噪声的平均功率（W）。这一公式深刻揭示了通信系统设计的两大核心追求：一是最大化信息传输速率，二是在给定信噪比条件下确保传输的可靠性。
围绕这一核心目标，通信系统在过去一百多年中经历了从模拟到数字、从窄带到宽带、从单天线到多天线、从单用户到大规模多用户的持续演进\cite{Goldsmith2005,Andrews2014}。从第一代移动通信系统（1G）的模拟蜂窝网络，到第五代移动通信系统（5G）的增强移动宽带（eMBB）、超可靠低时延通信（URLLC）和大规模机器类通信（mMTC）三大场景，通信系统的设计方法论始终围绕着信源编码、信道编码、调制解调、多址接入、资源调度等环节展开\cite{TS38300}。在通信系统中，无线信道被视为信息传输的“管道”，其时变特性（如多径衰落、多普勒扩展等）通常被视为需要克服或补偿的不利因素\cite{Goldsmith2005}。

几乎与无线通信技术同期发展，雷达（Radio Detection and Ranging, RADAR）技术自20世纪初萌芽，并在第二次世界大战期间获得了飞速发展\cite{Skolnik2008}。与通信系统形成鲜明对比的是，雷达/感知系统的核心目标不是传输信息，而是利用电磁波与环境、目标的交互作用，从回波信号中提取关于目标或环境的物理参数信息\cite{Skolnik2008,Richards2014}。具体而言，雷达系统的典型任务包括：
\begin{itemize}
    \item \textbf{目标检测（Detection）:} 
    判断特定空间区域内是否存在感兴趣的目标，这是一个经典的二元假设检验问题；
    \item \textbf{参数估计（Estimation）:}
    估计目标的距离、速度、角度等运动参数，其理论性能下界通常可由克拉美-罗下界（Cramér--Rao Lower Bound, CRLB）刻画\cite{Kay1993}；
    \item \textbf{目标识别与分类（Recognition/Classification）:} 
    通过分析回波信号的特征（如微多普勒特征）对目标类型进行判别；
    \item \textbf{环境成像与重建（Imaging）:}
    如合成孔径雷达（SAR）和逆合成孔径雷达（ISAR）通过高分辨率成像实现对目标或地形的精细描绘\cite{Richards2014}。
\end{itemize}
在雷达系统中，发射信号的设计追求良好的模糊函数（Ambiguity Function）特性——即在距离-多普勒二维平面上具有尖锐的主瓣和低旁瓣，以实现对目标参数的高分辨率估计\cite{Levanon2004,Richards2014}。这一设计理念与通信系统中追求的“频谱效率最大化”和“误比特率最小化”存在本质差异\cite{Liu2020JointRadarComm,Zhang2021ISAC6G}。
通信系统与雷达系统在核心目标上的根本差异，导致二者在频谱使用、硬件平台、信号设计和接收处理等关键技术维度上走上了截然不同的发展道路\cite{Liu2020JointRadarComm,Liu2022ISAC}。
\begin{enumerate}[label=(\arabic*)]
    \item \textbf{频谱使用}
    
    通信系统与雷达系统长期使用各自独立分配的频谱资源。在频谱管理的历史框架下，国际电信联盟（ITU）通过《Radio Regulations》以及世界无线电通信大会（WRC）为不同业务划分了专属频段。例如，蜂窝移动通信系统主要使用700 MHz至3.5 GHz等频段，而5G NR进一步扩展至毫米波频段；相比之下，各类雷达系统则分布在S波段（2--4 GHz）、C波段（4--8 GHz）、X波段（8--12 GHz）等不同频段。这种频谱的物理隔离使得两类系统可以在互不干扰的条件下独立运行。
    \item \textbf{硬件平台}
    
    通信设备（如基站、终端）与雷达设备在硬件架构上差异显著。通信系统的射频前端通常针对线性调制信号的发射/接收进行优化，强调信号链路的线性度与效率；而雷达系统，特别是脉冲雷达系统，其发射机常采用大功率功放，接收机则配置低噪声放大器和高动态范围的模数转换器（ADC），以应对回波信号微弱且动态范围极大的特殊需求。此外，雷达系统的天线设计也与通信系统存在显著差异——军用雷达通常采用高增益、窄波束的定向天线或相控阵天线，而蜂窝通信基站则需要兼顾覆盖范围的全向或扇区天线。
    \item \textbf{信号设计}
    
    通信信号的设计以承载调制信息为核心。典型的通信信号，如正交频分复用（OFDM）信号，通过在频域子载波上加载独立的数据符号来实现高效的信息传输。通信信号的波形设计关注的核心指标包括频谱效率、峰均功率比（PAPR）、带外泄漏等\cite{Goldsmith2005,Andrews2014}。相比之下，雷达信号的设计以优化探测性能为核心。经典的雷达信号包括线性调频（LFM）脉冲、相位编码脉冲等，其设计准则围绕模糊函数的形状展开——追求在距离维和速度维上同时获得高分辨率和低旁瓣\cite{Levanon2004}。雷达信号通常是确定性已知的，因为接收端需要利用已知的发射波形进行匹配滤波。而通信信号在雷达系统的视角下通常是“随机”的，因为数据符号对雷达接收来说并不是已知的先验信息\cite{Liu2020JointRadarComm,Zhang2021ISAC6G}。
    \item \textbf{接收处理}
    
    通信接收机的核心任务是信号检测与解码——即在存在噪声和干扰的条件下，从接收信号中恢复出发射端发送的信息比特。典型的处理流程包括同步、信道估计、均衡、解调与解码等步骤\cite{Goldsmith2005}。
    
    雷达接收机的核心任务则是目标检测与参数提取。典型的处理流程包括匹配滤波（脉冲压缩）、恒虚警率（CFAR）检测、多普勒处理、测角处理等。雷达接收机关注的是从回波信号中提取目标的物理参数，而非恢复调制信息\cite{Skolnik2008,Richards2014}。
\end{enumerate}
\begin{table}[!ht]
\centering
\begin{tabular}{|c|c|c|}
\hline
\textbf{对比维度} & \textbf{通信系统} & \textbf{雷达/感知系统} \\
\hline
\textbf{核心目标} & 可靠信息传输 & 目标检测与参数估计 \\
\hline
\textbf{频谱使用} & 授权通信频段 & 专用雷达频段 \\
\hline
\textbf{信号特性} & 随机调制数据承载 & 确定性已知波形 \\
\hline
\textbf{信道角色} & 信息传输的“管道”，需克服 & 感知的“目标载体”，需分析 \\
\hline
\textbf{发射信号设计准则} & 频谱效率、误码率 & 模糊函数、分辨率 \\
\hline
\textbf{接收处理目标} & 比特恢复 & 参数提取 \\
\hline
\textbf{性能度量} & 容量、吞吐率、延迟 & 检测概率、估计精度 \\
\hline
\textbf{典型硬件架构} & 基站/终端 & 雷达站/传感器 \\
\hline
\end{tabular}
\caption{通信系统与雷达/感知系统的核心对比}
\label{tab:comm-vs-radar}
\end{table}

表~\ref{tab:comm-vs-radar} 给出了传统通信系统和雷达感知系统在不同维度的主要区别。需要指出的是，通信系统与雷达系统的独立演进在5G及以前的网络世代中具有其历史合理性。首先，从频谱资源的角度看，在Sub-6 GHz频段中，通信系统与雷达系统所使用的频段存在较强的物理隔离，二者之间的频谱冲突问题尚不突出。其次，从业务需求的角度看，移动通信网络在5G及以前阶段主要服务于人与人之间的通信（语音、数据），网络并不需要对物理环境进行显式感知和理解。最后，从技术成熟度的角度看，通信与雷达各自的技术栈已经高度成熟且专业化，在各自领域内不断深化优化的路径是高效的。

然而，这种分立模式也正在暴露出日益严重的局限性。首先，频谱资源的紧张加剧了矛盾。随着移动通信向毫米波（mmWave）和太赫兹（THz）频段扩展，通信系统与雷达系统的频谱使用范围开始出现邻近甚至重叠，传统静态分配模式面临越来越大的压力\cite{TS38104,Liu2022ISAC}。例如，5G NR在24.25--52.6 GHz的毫米波频段与多种雷达应用（如汽车雷达、气象雷达）在频谱利用上存在显著邻近性\cite{Zheng2019RadarCommunicationCoexistence}。频谱资源的稀缺性使得两类系统长期独占各自频段的模式变得不可持续。
其次，新兴应用对“通信+感知”能力的联合需求不断涌现。自动驾驶车辆需要同时与路侧基础设施通信并对周围环境进行感知；无人机需要在飞行过程中同时进行遥控链路通信和避障感知；智慧城市中的海量物联网节点也需要在上传数据的同时感知周围环境\cite{Liu2022ISAC,De2021convergent}。在这些场景中，通信与感知功能的分立部署往往意味着更高的系统复杂度、成本和时延。
最后，独立部署导致的资源浪费日益显著。通信基站和雷达系统各自拥有射频前端、天线阵列、信号处理单元等硬件资源，但在很多场景下这些资源在物理层和平台层面具有共享潜力\cite{Liu2020JointRadarComm,Liu2022ISAC}。独立部署意味着硬件、能耗、站址等资源的重复投入。这些局限性共同构成了推动 ISAC 技术发展的根本动因，也为后续讨论的“通感融合需求”提供了历史背景。

\subsubsection{The fundamental demands for "convergence of communication and sensing" proposed by new wireless network scenarios}
\label{subsubsec:ch1-fundamental demands}

从根本上讲，ISAC（Integrated Sensing and Communications）的提出并非仅仅源于一种新的技术概念，而是新型无线网络应用场景对网络能力提出了本质性变化要求的结果。未来无线网络不再只需要提供传统意义上的“连接能力”，还必须进一步具备“环境感知能力”与“任务支撑能力”。在这一背景下，通信系统正在从单一的信息传输平台，逐步演化为兼具感知、认知与协同决策能力的综合信息基础设施\cite{Liu2022ISAC}。

首先，频谱资源紧张以及通信与感知双系统重复占用的问题，是推动通感融合的重要现实动因。频谱资源是无线通信系统与雷达/感知系统共同依赖的基础物理资源。随着全球移动数据流量持续增长，对频谱资源的需求不断攀升，而优质频谱尤其是低频和中频连续频段的稀缺性愈发突出。根据国际电信联盟（ITU）及各国频谱管理机构发布的频率划分与监管文件，在 Sub-6~GHz 频谱范围内，移动通信、卫星业务、无线电定位、航空航海雷达以及其他专用无线业务均已占据大量频谱资源，可供新系统独占使用的连续优质频段日益有限\cite{ITU_R_RR2024,miit_radio_freq_2026_en}。与此同时，移动通信系统正进一步向毫米波甚至太赫兹频段扩展，这使其不可避免地进入传统上由雷达、遥感或无线电定位系统长期使用或邻近使用的频率范围。以 24~GHz 附近频段为例，5G 毫米波业务已进入 24.25~GHz 以上频段，而 23.6--24.0~GHz 邻近频段长期用于地球探测卫星的被动遥感；此外，24~GHz 频段历史上也曾广泛用于短距汽车雷达。这表明，当通信系统不断向更高频段拓展时，其与传统感知/遥感系统之间的邻频共存甚至同频协调压力将显著增加。在此背景下，继续为通信与感知系统分别分配独立的频谱资源，不仅越来越困难，而且在总体频谱利用效率上也并不经济。ISAC 的一个核心优势就在于通过通信与感知共享频谱、共享基础设施和共享信号处理链路，从根本上缓解频谱资源紧张与双系统重复占用之间的矛盾；更进一步地，如果能够设计统一波形，使其同时承载通信信息并服务于环境感知，则系统所获得的收益将超越简单的分时或分频共享，从而带来额外的频谱效率增益\cite{Liu2022ISAC,Hassanien2015dual}。

其次，终端、基站与边缘节点对感知能力“内生化”的需求正在不断增强。传统网络体系中，感知通常由独立的传感器系统承担，例如雷达、摄像头或激光雷达等，而通信网络主要负责数据传输，两者之间的耦合程度相对较低。然而，面向 6G 的众多新型场景正在促使网络节点本身具备原生感知能力，即感知不再是附加于通信系统之外的外围模块，而逐步成为网络节点的一种基本功能\cite{Zhang2021ISAC6G}。这一趋势在多个典型场景中表现得尤为明显。例如，在车联网与自动驾驶场景中，车载通信模块已具备较强的射频、天线与基带处理能力，如果能够在通信链路上进一步实现目标检测、定位和环境感知，则有望减少对独立感知硬件的依赖，降低系统复杂度与部署成本；在低空经济与无人机网络场景中，飞行器往往需要在有限载荷和功耗约束下同时完成通信、定位、避障和环境感知任务，将感知能力嵌入通信链路有助于提升系统集成度与能效；在智能工厂和工业物联网场景中，海量设备不仅需要可靠连接，还需要实时获取人员、设备和物料的状态信息，如果由通信基础设施同时提供感知服务，则可显著减少额外传感器的大规模部署；在智慧城市和公共安全场景中，密集部署的基站若能够同时承担交通监测、异常事件检测与空间状态感知等任务，将大幅拓展现有网络基础设施的功能边界和服务价值。

更深层次地看，ISAC 的必要性来自于未来无线网络角色定位的根本性转变。在 5G 及以前的网络世代中，无线网络的核心任务主要是提供高效、可靠的信息传输，即网络更像是一种“信息传输管道”；网络关注的重点是速率、时延、覆盖和连接数，而对物理环境本身的理解和参与程度相对有限。然而，在面向 6G 的愿景中，无线网络正在从传统的“连接平台”逐步转变为“连接、感知、计算与控制”深度融合的综合服务平台\cite{Liu2022ISAC,Zhang2021ISAC6G}。这一转变意味着：网络不仅需要传输信息，还需要理解其所处的物理环境，例如感知用户的位置、速度与行为状态，刻画周围目标的分布特征，以及获取信道的空间结构信息；网络不仅需要提供链路，还需要支撑智能决策，例如在波束管理、资源调度、移动性管理和任务分配过程中，将通信状态信息与环境感知结果联合起来进行优化；对于自动驾驶、工业控制、远程操作等高可靠低时延闭环应用而言，网络还需要支持从环境感知、信息传输、边缘计算到执行控制的全流程协同，其中感知信息往往是形成闭环控制的起点。正因为如此，感知能力对于未来无线网络而言，不再是可有可无的附加组件，而正在成为网络原生能力的重要组成部分。

在理解上述需求之后，一个自然的问题是：既然未来网络需要感知能力，为什么不采用“通信系统 + 独立感知系统”的外挂式方案，而必须强调一体化设计？事实上，这种外挂式方案在当前工程实践中已较为常见，例如在道路基础设施中同时部署通信基站和独立雷达，在智能车辆中同时配置通信模组和毫米波雷达等。然而，随着应用场景对实时性、协同性、系统成本和规模部署能力提出更高要求，这种分离式架构的局限性也日益显现。其一，独立通信系统与感知系统之间往往需要额外的数据接口、同步机制和信息交互流程，这会带来不可忽略的处理开销与传输时延，在自动驾驶、工业控制等时延敏感场景中尤为不利。其二，通信与感知系统分别设计、各自优化，导致它们在信号体制、资源分配、空口协议和硬件平台等层面缺乏深度协同，难以实现全局最优的资源利用与性能平衡。其三，在每个节点同时部署通信设备和独立感知设备，意味着更高的硬件成本、功耗成本、站址成本与维护成本，这对于大规模、广覆盖的商用部署并不经济。其四，来自不同系统、不同制式的通信数据与感知数据在时间戳、坐标系、格式规范和处理流程等方面往往存在差异，后续融合处理复杂，系统实现难度较高\cite{Liu2022ISAC,Hassanien2015dual}。相比之下，ISAC 通过在系统架构、信号设计、资源调度和接收处理等层面进行一体化设计，能够从根本上提高频谱利用效率、增强感知与通信之间的协同能力，并为未来无线网络提供更高集成度、更低成本和更强任务支撑能力的统一技术路径。

\subsubsection{ISAC cannot be simply understood as a combination of "communications + radar"}
\label{subsubsec:ch1-isac-not-simple-combination}

明确了ISAC的必要性之后，还需要进一步澄清一个常见的概念性误解：ISAC 并不是将通信系统与雷达系统简单地“拼接”在一起。事实上，在已有文献与工程实践中，通信与感知的结合方式存在多种层次，从最基本的物理共址，到硬件共享、频谱共存，再到波形、协议与网络架构层面的深度融合，不同方案在“融合深度”上存在本质差异\cite{Liu2022ISAC,Zhang2021ISAC6G}。如果不对这些概念进行区分，容易将若干“通信与感知并存”的系统误认为 ISAC，从而模糊其真正的技术内涵。

最朴素的一种理解是：在同一个物理站点上同时安装一套通信设备和一套雷达设备，二者分别完成各自功能，彼此之间没有实质性的系统级协同。这种方式实现的只是空间上的共址（co-location），其主要优点在于部署方便，且有利于站址、供电和部分基础设施的复用。然而，从系统本质上看，它仍然是两套相互独立的设备并行部署，通信与感知在信号设计、资源调度、信息处理乃至性能目标上均相互分离，因此不能视为严格意义上的一体化系统。

比物理共址更进一步的一类方案，是在同一套硬件平台上复用部分基础组件，例如共享射频前端、天线阵列或基带处理平台，并通过时间切换、频率切换或工作模式切换的方式，在不同时隙、不同频点或不同配置下分别执行通信与感知功能。与完全独立部署相比，这类方案能够在设备体积、制造成本和能耗等方面带来一定优势，也体现出通信与感知功能之间的初步协同。然而需要指出的是，此时通信与感知任务在时间域或频率域上通常仍然是正交分离的，两种功能并未在统一信号、统一协议或统一资源优化框架下同时运行。因此，这类方案更适合被视为通信与感知的共平台实现，或 ISAC 的初级形态，而不能完全代表 ISAC 的完整内涵\cite{Hassanien2015dual}。

另一类容易与 ISAC 混淆的方案是频谱共存（spectrum coexistence）。在这类系统中，通信系统与雷达/感知系统共享同一段频谱资源，并通过认知接入、动态频谱分配、干扰抑制、波束隔离或干扰对齐等机制来降低彼此之间的相互影响\cite{Liu2020JointRadarComm}。频谱共存的研究重点在于：当两个原本独立设计的系统被迫或主动工作在同一频谱资源上时，如何保证二者能够“尽可能少地互相干扰”。因此，共存问题的核心是兼容性与干扰管理，而不是功能的统一设计与协同增强。换言之，在频谱共存框架下，通信和感知仍然是两个目标、架构和处理流程相互独立的系统，只不过共享了物理频谱资源；这与 ISAC 所追求的“在一套系统中同时实现通信与感知”的理念仍有明显差别\cite{Liu2022ISAC}。

从更严格的意义上说，真正的 ISAC 追求的是通信与感知功能在系统多个层面上的深层耦合（deep integration）。这种耦合并不局限于设备层面的共用，而是体现在信号、资源、硬件、协议、网络以及业务逻辑等多个层面上，其目标也不再只是“避免冲突”，而是通过联合设计实现通信性能与感知性能的协同增益\cite{Liu2022ISAC,Zhang2021ISAC6G}。具体而言，这种深层耦合主要体现在以下几个方面：

\begin{itemize}
    \item \textbf{信号层面的耦合：}同一信号波形同时承担通信信息承载与环境感知功能，而不是将通信信号与感知信号简单叠加，或通过交替发射分别完成两类任务。在理想情况下，通信数据本身构成感知波形设计的一部分，而感知回波中所反映的传播环境信息也能够反过来服务于通信链路估计、波束管理或信道预测等任务。

    \item \textbf{资源层面的耦合：}时域、频域、空域和功率域等无线资源由通信与感知任务共同占用，并在统一优化框架下进行分配与调度。此时，资源管理不再分别面向“通信系统”与“感知系统”独立进行，而需要同时考虑通信速率、时延、可靠性以及感知检测概率、估计精度和分辨率等多重性能指标之间的权衡。

    \item \textbf{硬件层面的耦合：}天线阵列、射频链路、模数/数模转换模块以及基带信号处理单元等底层硬件资源被通信与感知功能真正共享，其架构设计从一开始就同时面向两类任务需求，而不是简单地在同一硬件平台上轮流运行两套功能相互独立的软件或算法。

    \item \textbf{协议层面的耦合：}空口协议、帧结构、导频与参考信号设计、反馈机制、随机接入过程以及链路维护流程等，需要原生支持感知功能，而不是在传统通信协议之上附加一个感知功能模块。换言之，感知能力应当在协议设计之初就被纳入系统设计目标之中。

    \item \textbf{网络层面的耦合：}感知功能需要被纳入网络架构的整体设计，包括感知任务的触发与调度、感知数据的汇聚与分发、多节点之间的协同感知以及感知结果对网络资源管理的反向支持等。此时，感知不再只是单个节点的局部能力，而成为网络级能力的一部分。

    \item \textbf{业务层面的耦合：}通信服务与感知服务共同支撑上层应用，形成从“环境感知”到“信息传输”、再到“计算决策”和“控制执行”的完整业务闭环。对于自动驾驶、智能制造、低空协同和智慧城市等场景而言，这种业务层面的闭环正是 ISAC 价值的重要体现。
\end{itemize}

因此，ISAC 与传统“通信系统 + 雷达系统”的组合式方案之间最根本的区别，并不在于二者是否安装在同一个位置，也不在于是否共享部分硬件资源，而在于通信与感知是否在统一系统目标下实现了跨层联合设计与协同优化。只有当通信与感知不再被视为彼此独立的附属功能，而是在波形、资源、协议、网络与业务流程等层面形成有机统一时，系统才真正具有了通感一体化的本质特征。这种多层次、深耦合、面向任务闭环的融合机制，正是本书后续讨论的核心主线。

\subsection{The basic definition and conceptual boundaries of ISAC}
\label{subsec:ch1-basic-definition}
\subsubsection{The definition of ISAC}
\label{subsubsec:ch1-definition-isac}

在前文讨论了 ISAC 产生的现实动因、核心理念以及其与“简单共址”或“频谱共存”等方案之间的区别之后，本节进一步给出 ISAC 的基本定义，并对其概念边界作出较为清晰的界定。需要指出的是，尽管当前学术界与产业界对于 ISAC 的具体实现路径仍存在不同侧重，但对于其基本内涵已经形成较为一致的认识：ISAC 的核心不在于将通信功能与感知功能机械叠加，而在于在统一系统目标下实现通信与感知的联合设计、协同运行与双向赋能\cite{Liu2022ISAC,Liu2020JointRadarComm}。
本教材将 ISAC 定义如下：
\begin{quote}
\textbf{ISAC（Integrated Sensing and Communications，通感一体化）是指在统一的系统框架下，通过共享或协同设计的信号、硬件、频谱、空口与网络机制，使系统同时具备通信功能与感知功能的一类技术体系。}
\end{quote}
这一定义包含以下几个关键要素，需要分别加以说明。
\begin{enumerate}[label=(\arabic*)]
    \item \textbf{统一的系统框架}
    
    ISAC 的首要特征是“统一”，即通信与感知不是以两个彼此独立的子系统松散拼接而成，而是在一个共同的系统框架内进行建模、设计与运行\cite{Liu2022ISAC,Zhang2021ISAC6G}。这一统一框架通常体现在多个层面，包括统一的信号处理链路、统一的无线资源管理机制、统一的网络架构以及统一的性能评估视角等。需要强调的是，这里的“统一”并不意味着通信与感知在物理目标、性能指标或处理方法上完全相同。相反，二者在速率、时延、可靠性、检测概率、估计精度、分辨率等方面往往具有不同甚至相互冲突的需求。所谓统一，强调的是这些差异不是彼此割裂、分别优化，而是在系统设计之初即被纳入同一个整体框架中进行权衡与协调。
    \item \textbf{共享或协同设计}
    
    在 ISAC 系统中，通信与感知的融合关系通常体现为“共享”与“协同设计”两种基本形式，而实际系统往往处于这两者之间的连续谱上\cite{Liu2020JointRadarComm,Gonzalez2024integrated}. 一方面，在共享模式下，通信与感知功能可以共用相同的信号波形、频谱资源、硬件平台或处理模块。例如，同一 OFDM 信号既可用于向用户传输信息，也可通过分析其传播路径或回波结构来提取环境感知信息。另一方面，在协同设计模式下，通信与感知未必完全使用同一组资源或完全相同的信号形式，但二者在系统设计阶段是联合考虑、共同优化的。例如，在波束赋形或功率分配过程中，需要同时考虑通信覆盖、数据速率与感知方向性、探测精度之间的平衡，从而实现整体系统性能的最优折中。由此可见，ISAC 并不局限于“所有资源完全共用”的极端情形，更一般地，它强调的是通信与感知在设计目标和资源使用上的内在协同性。
    \item \textbf{同时具备通信功能与感知功能}
    
    ISAC 定义中的“同时”应从系统目标和运行机制两个层面来理解。首先，在系统目标层面，通信与感知在设计之初就是两个并列的核心功能，而不是“以通信为主、附带一点感知”或“以感知为主、顺便传一点信息”的从属关系\cite{Liu2022ISAC}. 其次，在系统运行层面，“同时”既可以表现为严格意义上的并行工作，即同一时刻、同一信号或同一资源块同时承担通信与感知任务；也可以表现为在极短时间尺度上的快速切换与准同时执行。对于后一种情况，尽管从瞬时角度看通信与感知可能并非完全重合，但由于二者在系统逻辑、资源控制和任务调度上已经实现统一管理，因此仍属于 ISAC 的范畴。换言之，判定一个系统是否属于 ISAC，关键不在于通信与感知是否绝对同步发生，而在于两种功能是否在统一系统框架下被对等建模、联合设计与协同运行。
    \item \textbf{感知对象的广泛性}
    
    在 ISAC 语境下，“感知”应被理解为一个广义概念，其对象远不限于传统雷达中典型的点目标或飞行器目标\cite{Liu2022ISAC,Gonzalez2024integrated}. 更一般地，ISAC 所面向的感知对象至少可以包括以下几类：第一，外部物理环境，例如建筑物、墙体、道路结构、地形甚至天气条件等静态或准静态环境要素；第二，运动目标，例如车辆、行人、无人机等动态目标的检测、定位、跟踪与识别；第三，网络对象，例如通信用户的位置、移动状态、空间分布与行为模式；第四，设备与终端，例如物联网设备的存在性检测、状态估计和活动监测；第五，空间电磁环境，例如传播路径、散射体分布、角度-时延结构以及干扰状态等；第六，生命体征与人体活动，例如呼吸监测、心跳估计、跌倒检测、手势识别以及人体行为识别等。正是由于感知对象具有如此广泛的内涵，ISAC 的应用边界远远超出了传统雷达系统的范畴，也使其与未来无线网络、环境智能和人机交互等领域形成了更加紧密的联系。
    \item \textbf{通信与感知的双向赋能关系}
    
    ISAC 的另一个本质特征在于，通信与感知之间并非单向依附关系，而是能够形成双向赋能\cite{Liu2022ISAC,Zhang2021ISAC6G}. 一方面，感知可以辅助通信，即通过获取用户位置、目标运动状态、障碍物分布、空间散射结构或链路可视性等信息，为通信系统的波束管理、链路建立、信道预测、切换决策和资源调度提供先验信息，从而提高通信链路的可靠性、稳定性和资源利用效率。另一方面，通信也能够赋能感知，即通信系统天然具备射频信号源、网络覆盖、时间同步、分布式节点部署和反馈机制等基础条件，这些条件可为环境感知、目标检测与多节点协同感知提供天然载体和基础设施支持。特别是在蜂窝网络、Wi-Fi 网络和车联网等大规模部署场景中，通信网络的多节点分布和持续在线特性，使其相较于传统独立雷达更适合作为广域感知平台。正是这种“感知辅助通信”与“通信赋能感知”的双向作用机制，使得 ISAC 不应被理解为“在通信系统中附带加入一个感知模块”，而应被理解为一种通过双功能融合实现系统整体性能跃升的统一技术框架。
\end{enumerate}
综合来看，ISAC 的基本定义至少包含三个层面的核心含义：其一，通信与感知在系统层面是统一设计、统一运行的；其二，通信与感知之间不仅共享资源，而且能够通过联合优化形成协同增益；其三，感知在 ISAC 中是广义的、网络化的、能够反向服务于通信系统优化的。对这些关键要素的把握，有助于我们在后续讨论具体信号模型、系统架构和性能权衡问题时，更准确地理解 ISAC 的研究对象与技术边界\cite{Liu2022ISAC,Liu2020JointRadarComm,Gonzalez2024integrated}。

\subsubsection{The boundary distinction between ISAC and related concepts}
\label{subsubsec:ch1-boundary-distinction}

在 ISAC 的学术研究、工程实践与标准化讨论中，存在多个与之密切相关但又不完全等同的技术概念。准确区分 ISAC 与这些概念之间的联系与差异，对于建立清晰的技术图景、避免术语混用以及理解不同研究工作的实际定位都十分重要。总体而言，这些相关概念大致可以分为三类：第一类是以干扰管理为核心的“共存”问题；第二类是以双功能波形和链路级联合设计为核心的“联合雷达通信”或“双功能雷达通信”问题；第三类是术语层面与 ISAC 相近、但强调重点略有差异的相关表述，如 JCAS、JRC 和 RadCom 等。

\paragraph{1) ISAC 与 RCC 的区别}

雷达--通信共存(Radar-Communication Coexistence)是通信与雷达结合问题中最早受到广泛关注的研究方向之一。其核心问题是：当原本独立设计的通信系统和雷达系统由于频谱紧张、邻频部署或同频工作等原因被迫共享同一频段或相邻频段时，如何通过合理的干扰建模、资源协调和信号处理机制，使二者能够在可接受的性能损失下共同运行\cite{Chiriyath2017RadarCommCoexistence,Liu2020JointRadarComm}。围绕这一问题，典型研究内容包括雷达信号对通信接收机的干扰分析与抑制、通信信号对雷达检测与参数估计性能的影响评估，以及基于功率控制、波束零陷、认知频谱接入、频率规划和空时干扰抑制等手段的共存策略设计\cite{Deng2013InterferenceMitigation,Sodagari2012ProjectionBasedCoexistence,Zheng2019RadarCommunicationCoexistence}。

从本质上看，雷达--通信共存关注的是\textbf{两个独立系统之间的兼容性问题}。在这种框架下，通信系统和雷达系统分别服务于各自独立的任务目标，彼此之间的主要关系是“避免互相破坏”或“在冲突中尽量协调”。因此，共存研究的出发点是冲突管理，其关键目标是降低相互干扰并保证各自功能尽可能正常运行。与之相比，ISAC 关注的则是\textbf{在统一系统框架下同时实现通信与感知功能}，其目标不是简单减少两个系统之间的摩擦，而是在同一系统目标下实现通信与感知的联合设计、协同优化与性能增益\cite{Liu2022ISAC}. 换言之，共存强调“两个系统如何和平共处”，而 ISAC 更强调“一个系统如何原生地同时具备两种能力”。这一差异是理解 ISAC 概念边界的首要前提。

\paragraph{2) ISAC 与 DFRC/JRC 的关系}

双功能雷达--通信（Dual-Function Radar-Communication, DFRC）系统通常被视为 ISAC 的重要技术来源之一，也是通信与感知融合研究中最具代表性的分支方向之一\cite{Hassanien2015dual,Liu2020JointRadarComm}。DFRC 的核心思想在于设计一种能够同时承担信息传输与目标探测功能的双功能信号或发射机制，使得同一发射波形既可被通信接收机用于解调信息，又可由雷达接收机通过回波处理提取目标的距离、速度、角度等参数。围绕 DFRC，研究者已经提出了大量具有代表性的物理层方法，例如基于 OFDM 的双功能波形设计、联合波束赋形、嵌入式信息调制、符号级波形设计以及 MIMO 雷达通信融合架构等\cite{Hassanien2015dual,Zhang2021ISAC6G}。

从关系上看，DFRC 可以被理解为 ISAC 在\textbf{物理层信号设计与链路级收发机设计}维度上的重要实现形式。也就是说，DFRC 所研究的许多问题，如双功能波形设计、联合发射策略和通信/雷达性能权衡等，构成了 ISAC 技术体系中的关键组成部分。然而，ISAC 的概念外延通常比 DFRC 更广\cite{Liu2022ISAC,Gonzalez2024integrated}。DFRC 往往更侧重于“如何利用一个信号或一个发射平台同时服务于雷达和通信”，其分析重点主要集中在物理层、波形层、阵列处理层以及链路级性能上；而 ISAC 除了包含这些内容之外，还进一步强调网络架构、协议机制、资源调度、多节点协同感知、感知信息回传与融合处理、业务闭环构建等更高层问题。因此，可以说 DFRC 是 ISAC 的一个重要子集或关键实现方向，但 ISAC 并不等同于 DFRC；前者是一种更广义的系统级技术框架，后者则更偏向其中的链路级和波形级研究范式。

需要进一步说明的是，JRC（Joint Radar and Communication）一词在文献中与 DFRC 存在较大重叠，二者往往都用于描述雷达与通信功能的联合设计问题\cite{Liu2020JointRadarComm,Zhang2021ISAC6G}。在不少场合，JRC 更强调“联合设计”的思想，而 DFRC 更强调“单一平台或单一波形承担双重功能”的实现特征；不过，这种区分并非绝对，不同文献中的用法也不完全一致。因此，在阅读相关研究时，应结合具体上下文判断其所指范围，而不宜仅凭术语字面作出机械划分.

\paragraph{3) ISAC 与 JCAS、JRC、RadCom 等术语的区别}

除 DFRC 之外，文献中还常见若干与 ISAC 含义接近的术语，如 JCAS（Joint Communication and Sensing）、JRC（Joint Radar and Communication）以及 RadCom（Radar-Communication）等\cite{Liu2022ISAC,Gonzalez2024integrated}。这些术语之间既存在交叉，也反映出研究者在关注重点和学术传统上的差异。

其中，JCAS（联合通信与感知）与 ISAC 的语义最为接近，在很多论文中二者甚至可以近似互换使用。相较而言，JCAS 更突出“joint”，即通信与感知在设计上的联合性；而 ISAC 中的“integrated”则更强调二者在系统架构、资源管理、协议机制与业务流程等层面上的深层融合。因此，从语义细微差别上看，ISAC 通常比 JCAS 更强调系统级的一体化特征，但在具体文献中，两者往往代表相近甚至重叠的研究对象。

JRC 则更明确地将“雷达”作为感知功能的主要代表，因此它在传统雷达和通信融合研究中使用较多，特别适合描述以雷达探测、参数估计和目标跟踪为核心任务的系统\cite{Liu2020JointRadarComm,Hassanien2015dual}。相比之下，ISAC 使用更广义的 “sensing” 概念，其感知对象不仅包括传统雷达意义下的目标，也可能包括用户、设备、环境、电磁传播结构乃至人体活动等。因此，JRC 在概念上通常可以被视为 ISAC 的一个更具体、更偏雷达语境的子集。

RadCom 或 Radar-Communication 通常是一种较宽泛的领域性称谓，可用于泛指雷达与通信结合相关的研究方向，其边界往往依赖于具体文献语境。有些文献用其表示共存问题，有些则用其表示联合设计，因而其定义并不总是严格统一。正因如此，在教材中若不加说明地直接使用这些术语，容易给初学者带来混淆。

基于上述考虑，本教材统一采用 ISAC（Integrated Sensing and Communications）作为核心术语。这一选择主要基于以下三点原因：其一，“Integrated”能够较准确地体现通信与感知在系统设计中的深层耦合与一体化特征；其二，“Sensing”相较于“Radar”具有更强的广义性，能够覆盖雷达探测、定位、环境感知、人体活动识别以及信道环境感知等多种形式；其三，ISAC 已经成为当前学术界与产业界讨论通信与感知融合问题时最广泛使用的重要术语之一，并在 6G 愿景研究与相关标准化讨论中受到持续关注\cite{Liu2022ISAC,Gonzalez2024integrated,De2021convergent}。

\paragraph{4) “Integrated” 的多层含义}

最后，还需要进一步强调的是，ISAC 中的“Integrated”并不是一个单一维度上的概念，而是一个具有层次性的系统性概念。所谓“一体化”，并不只是指通信与感知共享某一段频谱，或共用某一套硬件，而是可能体现在从物理层到业务层的多个层面上\cite{Liu2022ISAC,Zhang2021ISAC6G,Liu2020JointRadarComm}。从教学和系统分析的角度出发，可以将这种一体化大致概括为以下几个层面：

\begin{itemize}
    \item \textbf{信号层一体化：}通信与感知使用统一波形，或在波形设计中同时面向信息承载与参数感知两类目标进行联合优化。

    \item \textbf{硬件层一体化：}天线阵列、射频链路、模数/数模转换器以及基带处理平台等底层硬件由通信与感知共同使用，并在架构设计上同时满足两类功能需求。

    \item \textbf{算法层一体化：}信号处理与优化算法同时服务于通信与感知目标，例如联合信道估计与目标检测、联合波束赋形、联合资源调度与状态感知等。

    \item \textbf{网络层一体化：}感知功能被纳入网络架构整体设计，包括感知任务的触发与调度、感知数据的上传与分发、多节点协同感知以及感知结果对通信网络优化的反馈支持。

    \item \textbf{业务层一体化：}通信服务与感知服务面向上层任务进行统一编排，形成“感知--通信--计算--决策--控制”的闭环业务流程，以支撑自动驾驶、工业控制、低空协同和智慧城市等复杂应用场景。
\end{itemize}

需要指出的是，不同 ISAC 系统在上述各层面上的融合深度可能并不相同，因此现实中的 ISAC 并不是一个“非此即彼”的概念，而更适合被理解为一个从浅层融合到深层融合的连续谱系。也正因如此，在评价某一系统是否属于 ISAC 时，不能仅凭其是否共享部分频谱或硬件作出简单判断，而应综合考察其在信号、资源、协议、网络与业务流程等多个层面上的融合程度与联合设计深度。

\subsubsection{The working definition of ISAC in this textbook}
\label{subsubsec:ch1-working-definition}

在前文对 ISAC 的基本定义及其与相关概念的边界进行了讨论之后，本教材还需要进一步给出一个面向教学、建模与后续分析的\emph{工作定义}（working definition）。之所以有必要引入工作定义，是因为 ISAC 作为一个仍在快速发展的研究方向，其具体实现形态、系统边界与应用重点在不同文献中并不完全一致。对于教材而言，仅给出一般性定义还不够，还必须明确本书讨论问题的基本立场、抽象层次与分析范围，以保证后续章节在理论模型、系统架构和性能指标上的一致性。
基于上述考虑，本教材将采用如下 ISAC 理解框架。
\begin{enumerate}[label=(\arabic*)]
    \item \textbf{将 ISAC 视为 6G 及未来无线网络的重要原生能力之一:}
    本教材将 ISAC 置于 6G 及未来无线网络的总体技术框架下进行讨论。与传统网络主要强调连接能力不同，面向未来的无线网络被普遍认为需要同时具备通信、感知、计算与控制等多维能力，其中感知不再仅仅依赖外挂式独立传感器系统，而正在逐步演化为网络节点和网络基础设施的原生能力之一\cite{Liu2022ISAC,Gonzalez2024integrated,Saad2020Vision6G}. 在这一定位下，本教材对于 ISAC 的讨论将不局限于点到点链路上的物理层信号设计，而是同时覆盖链路层、网络层乃至系统级架构问题，包括蜂窝组网、多节点协同感知、边缘计算支撑以及端--边--云协同处理等内容。换言之，本教材所讨论的 ISAC，不是狭义的“某种双功能波形设计问题”，而是面向未来无线网络演进的一类系统级能力框架。
    \item \textbf{聚焦通信与感知功能协同设计的系统性问题:}
    本教材的核心关注点，不是孤立地讨论某一种波形、某一个算法或某一种应用，而是关注通信功能与感知功能如何在统一系统中进行协同设计，以实现整体性能的优化与任务支撑能力的提升\cite{Liu2020JointRadarComm,Zhang2021ISAC6G}. 具体而言，本书后续将围绕以下几类问题展开：
    \begin{itemize}
        \item 如何设计能够同时服务于通信与感知的信号模型与发射波形？
        \item 如何在通信与感知之间共享、分配并联合优化时频空功率等无线资源？
        \item 如何在 MIMO、多天线和大规模阵列系统中设计兼顾通信质量与感知性能的联合波束赋形策略?
        \item 如何建立统一的性能分析框架，以刻画通信速率、误码率、时延与感知检测概率、估计精度、分辨率之间的内在权衡？
        \item 如何在网络架构层面支撑大规模、多节点、分布式和协同式通感一体化服务？
    \end{itemize}
    上述问题分别对应 ISAC 研究中的信号层、资源层、阵列处理层、性能评估层和网络架构层，它们共同构成了本教材的主要技术主线。
    \item \textbf{不预设唯一最优波形、唯一最优架构或唯一应用场景:}
    ISAC 是一个具有明显开放性与多样性的技术体系，并不存在放之四海而皆准的唯一最优波形、唯一最优系统架构或唯一标准应用范式\cite{Liu2022ISAC,Liu2020JointRadarComm}. 不同频段、不同节点能力、不同网络部署方式以及不同业务目标，都会导致 ISAC 系统在设计上呈现出显著差异。因此，本教材不将 ISAC 简化为某一种固定实现方案，而是尽可能从统一原理出发介绍多种具有代表性的技术路径，包括但不限于：基于 OFDM 的通感一体化波形设计、基于 MIMO 的联合波束赋形方案、蜂窝网络中的感知组网与协同处理机制、毫米波/太赫兹频段的高分辨感知通信系统，以及面向车联网、低空经济、智能工厂和智慧空间等典型场景的 ISAC 应用框架\cite{Liu2022ISAC,Gonzalez2024integrated}. 通过这种方式，读者能够理解 ISAC 的共性思想，同时认识其在不同技术条件与应用场景下的差异化实现形式。
\end{enumerate}
上述工作框架不仅规定了本教材的讨论范围，也为后续章节建立统一的数学建模与性能分析语言奠定了基础。具体而言，在本书后续的理论分析中，通信性能通常将主要通过信道容量、可达速率、误比特率、可靠性和传输时延等指标进行描述；感知性能则主要通过检测概率、虚警概率、距离/速度/角度分辨率、估计误差以及克拉美--罗下界（Cram\'er--Rao lower bound, CRLB）等指标进行刻画\cite{Liu2022ISAC,Zhang2021ISAC6G}. 当通信与感知需要在统一系统中协同工作时，两类性能指标之间往往并不存在简单的单调一致关系，而需要通过多目标优化、约束优化或帕累托最优（Pareto optimality）等分析工具来刻画其折中与协同边界\cite{Liu2020JointRadarComm}. 在系统建模方面，本教材将尽可能采用统一的数学语言，对发射信号模型、传播信道模型、目标/散射体模型、观测模型和接收处理流程进行规范化描述，并在此基础上逐步建立联合性能评估框架，用于分析速率、可靠性、时延、定位精度、检测概率、估计误差和能耗等指标之间的关系。
需要特别指出的是，这种工作定义并不意味着本教材将 ISAC 限定为某一种狭义技术实现，而是意味着：在后续章节中，只要某种系统满足“在统一系统框架下，通过通信与感知功能在信号、资源、硬件、协议与网络机制上的共享或协同设计，实现信息传输与环境/状态感知的联合服务，并体现二者之间的双向赋能关系”这一基本特征，便将其纳入本教材的 ISAC 讨论范围\cite{Liu2022ISAC,Gonzalez2024integrated}. 相反，对于那些仅仅实现物理共址、简单频谱共存或功能松散拼接、但缺乏统一系统目标与联合设计机制的方案，本教材一般不将其视为严格意义上的 ISAC 核心范畴。
基于以上约定，本教材采用如下工作定义：
\begin{quote}
\textbf{ISAC 是面向 6G 及未来无线网络的一种关键原生能力，指在统一系统框架下，通过通信与感知功能在信号、资源、硬件、协议与网络机制上的共享或协同设计，实现信息传输与环境/状态感知的联合服务，并通过二者的双向赋能支撑网络智能化、场景理解与业务闭环。}
\end{quote}
这一工作定义将作为本书后续各章展开系统模型建立、算法设计、性能分析和应用讨论的共同出发点。

\subsection{The three-stage logic of ISAC evolution}
\label{subsec:ch1-evolution-logic}

ISAC 技术的形成并非一蹴而就，而是经历了一个从“相互独立但被迫共存”逐步演进到“统一框架下深度融合”的过程。为了帮助读者把握这一发展脉络，本教材将通信与感知融合技术的演进逻辑概括为三个具有明显递进关系的阶段：\emph{频谱共存与交叉应用阶段}、\emph{双功能平台与联合信号设计阶段}以及\emph{面向 6G 的内生通感一体化阶段}。这三个阶段并不是严格割裂的历史片段，而更适合被理解为研究重点不断深化、系统融合程度不断增强的技术谱系。理解每一阶段的核心目标、典型方案与内在局限，有助于把握 ISAC 从“共存”走向“融合”的内在逻辑。

\subsubsection{Stage 1: Spectrum Coexistence and Cross-Application Phase}
\label{subsubsec:ch1-stage1}

ISAC 演进的第一阶段，可以概括为频谱共存与交叉应用阶段。该阶段大致起源于通信系统与雷达系统在频谱资源上开始出现显著冲突的时期，其核心问题不是“如何将通信与感知做成一个系统”，而是“当两个原本独立的系统必须共享或邻接使用同一频谱资源时，如何尽量降低彼此干扰并维持各自功能的正常运行”\cite{Liu2020JointRadarComm,Chiriyath2017RadarCommCoexistence}。因此，这一阶段的基本特征在于：通信系统与雷达系统在目标、架构、波形、协议和性能评价体系上仍然彼此独立，但在频谱层面需要进行一定程度的协调和兼容。

从系统视角看，这一阶段的核心矛盾来自频谱资源紧张与双系统独立部署之间的冲突。通信系统和雷达系统各自采用独立的信号格式、硬件平台、接收处理流程和优化目标：前者通常追求更高的传输速率、更低的误码率和更好的链路可靠性，后者则更关注目标检测概率、虚警概率、参数估计精度与空间分辨能力。在这种背景下，研究与工程实践的重点集中于如何构建兼容机制，而不是构建统一系统\cite{ITU_R_RR2024}。换言之，这一阶段的联合仅发生在“频谱使用协调”层面，而尚未延伸到信号设计、资源优化和网络架构的一体化层面。
围绕这一阶段，典型研究问题主要包括以下几个方面：
\begin{itemize}
    \item \textbf{频谱冲突分析：}定量评估通信系统与雷达系统在同频或邻频工作条件下的相互影响，包括互干扰信号的统计建模、干扰功率计算、接收信干噪比分析以及干扰对通信误码率、雷达检测性能和参数估计精度的影响评估\cite{Chiriyath2017RadarCommCoexistence}。
    \item \textbf{干扰管理：}设计降低相互干扰的技术机制，例如发射功率控制、方向性波束隔离、保护频带设计、动态频谱协调、空时滤波以及接收端的干扰抵消等\cite{Liu2020JointRadarComm}。
    \item \textbf{共存机制设计：}建立通信系统与雷达系统在共享频谱条件下的工作规则，例如基于感知与认知的动态频谱接入协议、频谱占用状态检测、共享时序协调以及面向监管要求的频谱复用策略等\cite{Chiriyath2017RadarCommCoexistence}.
\end{itemize}
从本质上说，共存阶段中的通信与雷达系统仍然遵循“各自独立优化、相互兼容运行”的思想。通信系统通常以传输效率和链路质量为中心进行设计，雷达系统则以探测与估计性能为核心目标；两个系统之间并不存在真正意义上的统一目标函数，也缺乏贯穿发射、接收、资源调度和协议控制全过程的联合优化机制。因此，这一阶段的技术边界可以概括为：\textbf{在干扰管理层面存在协调，但在系统设计层面尚无真正融合}。这也是其与后续 ISAC 阶段的根本区别。

在具体实现上，共存阶段出现了若干具有代表性的技术方案。这些方案虽然尚未达到一体化设计的程度，但为后续通信与感知的深层融合奠定了问题基础和方法基础。
\begin{itemize}
    \item \textbf{时分共享（Time-Division Sharing）：}通信系统与雷达系统在时间维度上交替使用同一频段，例如在雷达脉冲间隙安排通信传输，或按照预定义时隙分别执行感知与通信任务。这种方案实现机制相对简单，系统间隔离度高，便于满足监管与工程部署要求，但其代价是频谱资源无法被真正同时利用，整体频谱效率提升有限\cite{Chiriyath2017RadarCommCoexistence}。
    \item \textbf{频分隔离（Frequency-Division Isolation）：}在共享或邻近频段内，将通信系统与雷达系统分配至不同子频带，并通过设置保护频带、滤波抑制和频率规划等方式降低邻道干扰。这类方案在工程上较为直观，也易于与现有频谱管理框架兼容，但本质上仍然属于资源切分而非功能融合，难以充分发挥宽带频谱资源的潜力\cite{ITU_R_RR2024}。
    \item \textbf{共存感知或干扰感知处理（Coexistence-Aware Sensing/Communication）：}在接收端或信号处理端显式考虑来自另一系统的干扰结构。例如，雷达接收机将通信信号建模为结构化干扰并进行抑制，或者通信接收机对雷达脉冲干扰进行检测、估计与消除。这类方法的特点在于不改变两系统的独立性，而是通过更智能的接收处理来提升共存能力\cite{Chiriyath2017RadarCommCoexistence,Liu2020JointRadarComm}。
    \item \textbf{机会式频谱复用（Opportunistic Spectrum Reuse）：}借鉴认知无线电思想，系统通过频谱感知、占用检测或环境认知判断某一频段是否处于空闲或低干扰状态，并在满足约束的条件下机会性复用该频谱资源。这类方法在理论上能够提高频谱利用率，但对频谱状态感知准确性、快速接入机制和监管兼容性提出了较高要求\cite{Chiriyath2017RadarCommCoexistence,Haykin2005CognitiveRadio}.
\end{itemize}

尽管上述共存方案在一定程度上缓解了通信系统与雷达系统之间的频谱冲突，但其局限性也十分明显。首先，在共存框架下，一个系统性能的改善往往需要以牺牲另一个系统的性能为代价，因而整体系统性能提升空间有限；其次，由于通信系统与雷达系统在硬件、波形、协议和信息处理链路上仍保持独立，难以形成统一架构，也难以实现真正意义上的深度资源共享；再次，两类系统之间缺乏实质性的双向信息交互与联合优化机制，因而难以释放感知辅助通信、通信赋能感知所可能带来的协同增益；最后，从工程实现角度看，维护两套独立系统及其频谱协调机制，本身也会带来额外的复杂度与部署成本\cite{Liu2020JointRadarComm,Liu2022ISAC}. 正是这些内在局限，推动了技术研究从“如何共存”进一步迈向“如何在同一平台和同一信号层面实现双功能”，并由此进入 ISAC 演进的第二阶段。

\subsubsection{Stage 2: Dual-Functional Platform and Joint Signal Design Phase}
\label{subsubsec:ch1-stage2}

如果说第一阶段的关键词是“共存”，那么第二阶段的关键词就是“联合设计”。这一阶段标志着通信与感知关系从“两个独立系统之间的兼容协调”进一步发展为“同一平台或同一信号层面的双功能实现”，从而构成了 ISAC 从共存走向融合的第一次实质性跃迁。该阶段最具代表性的思想可以概括为“一个信号、两种用途”（one signal, dual purposes）或“一个平台、双重功能”（one platform, dual functions）：通信与感知不再分别依附于两套彼此独立的发射机、接收机和波形体系，而是开始共享同一发射波形、同一天线阵列、同一射频前端，甚至共享部分接收处理链路。

从技术目标看，第二阶段的核心问题已经不再是单纯减少两个系统之间的干扰，而是如何通过联合设计使同一信号或同一平台同时满足通信与感知的双重性能需求。也正是在这一阶段，双功能雷达--通信（DFRC）、联合雷达通信（JRC）以及早期 ISAC 物理层方法逐步形成了较为完整的技术体系\cite{Liu2020JointRadarComm,Zhang2021ISAC6G,Hassanien2015dual}。围绕这一目标，典型研究主要集中在以下几个方面：
\begin{itemize}
    \item \textbf{双功能波形设计（Dual-Function Waveform Design）：}如何设计一个同时适用于数据传输与目标感知的发射波形，使其既具有足够的信息承载能力，又具备适合目标检测、参数估计和距离--速度解析的时频特性\cite{Liu2020JointRadarComm,Hassanien2015dual}。这一问题的根本挑战在于，通信与感知对波形性质的要求往往并不完全一致：通信通常希望波形具有较强的信息调制自由度、良好的频谱效率和抗多径能力，而感知则更关注波形的自相关特性、模糊函数形状、带宽利用方式以及对目标参数的可辨识性。因此，双功能波形设计本质上是一个在通信可达速率与感知估计性能之间寻求平衡的联合优化问题。
    \item \textbf{联合波束赋形（Joint Beamforming）：}在多天线与 MIMO 系统中，如何设计发射预编码矩阵或波束赋形向量，使其既能够为通信用户提供足够的接收信噪干比（SINR）或速率保障，又能够在目标方向形成所需的照射功率、空间聚焦能力或角度分辨能力\cite{Liu2020JointRadarComm,Liu2022ISAC}。这一问题通常具有明显的多目标优化特征，因为通信波束往往偏向用户方向，而感知波束则需要覆盖目标区域或形成特定的波束图结构，二者在空间资源分配上可能存在冲突。
    \item \textbf{共享资源调度与联合优化（Shared Resource Scheduling and Joint Optimization）：}当通信与感知共享同一组时间、频率、空间和功率资源时，如何在给定系统约束下实现资源的最优分配，是第二阶段的另一核心问题\cite{Zhang2021ISAC6G,Liu2022ISAC}。这一问题既可以体现在单帧内的导频/数据/感知功能分配，也可以体现在子载波选择、功率加载、阵列自由度分配和时隙结构配置等层面。本质上，这类问题需要在通信性能指标与感知性能指标之间建立统一的折中分析框架。
\end{itemize}
尽管第二阶段已经在物理层与链路层实现了通信与感知的实质性联合设计，但其研究边界总体上仍然主要集中于单链路、单节点或单小区场景。具体而言，联合波形设计通常围绕点到点或点到多点链路展开，联合波束赋形多在单基站或单发射平台作用范围内进行分析，而性能评估也主要基于链路级理论分析或仿真验证\cite{Liu2020JointRadarComm,Liu2022ISAC}. 相比之下，多小区协同感知、网络级感知任务调度、感知结果的汇聚与融合、感知服务的端到端编排以及感知与边缘计算协同等问题，在这一阶段尚未成为研究重点。因此，从技术谱系的角度看，第二阶段虽然奠定了 ISAC 的核心理论基础，但其系统视野仍然偏向链路级和平台级，尚未完全拓展到网络级原生能力框架。

在具体技术实现方面，第二阶段形成了若干具有代表性的方案类型。这些方案构成了后续 ISAC 研究最重要的技术基础。
\begin{itemize}
    \item \textbf{基于 OFDM 的通感一体化方案（OFDM-Based ISAC）：}OFDM 是第二阶段中最具代表性的双功能信号载体之一，因为它天然具有规则的时频结构、成熟的通信收发机制以及较强的参数处理便利性\cite{Liu2022ISAC,Sturm2011WaveformDesign,Zhou2022integrated}. 在这类方案中，OFDM 信号既用于承载通信数据，也可通过对回波信号在子载波维和符号维上的联合处理来提取目标的距离与速度信息。由于 OFDM 的二维时频结构与雷达中的距离--多普勒处理具有天然对应关系，许多方法可通过二维快速傅里叶变换（2D-FFT）或其变形实现较为高效的参数估计。这类方案的重要意义在于，它直接验证了主流通信波形可以在较小改动下承担基本感知任务，从而极大促进了通信系统向 ISAC 平台的演进。
    \item \textbf{DFRC 波形设计：}另一条重要技术路线是在传统雷达波形中嵌入通信信息，或在通信信号中显式引入面向雷达感知的结构约束，从而形成真正意义上的双功能波形\cite{Hassanien2015dual,Blunt2016overview,Blunt2010Embedding}。这类方法包括但不限于：在脉冲雷达或线性调频（LFM）信号中嵌入调制信息、通过天线索引或波束方向切换实现空间调制式信息传输、在满足通信质量约束条件下优化发射波形的模糊函数、以及在 MIMO 雷达波形中联合设计信息嵌入机制等。相较于直接复用现有通信波形，DFRC 方案往往更强调从雷达性能出发构造双功能信号，因此在目标检测与角度控制方面具有一定优势，但也常常面临通信兼容性、标准适配性和系统复杂度方面的挑战。
    \item \textbf{MIMO 联合设计与波束赋形：}随着多天线技术的发展，MIMO 成为第二阶段最核心的联合设计平台之一\cite{Liu2020JointRadarComm,Li2023integrated}. 在这类方案中，发射端通过联合设计预编码矩阵、协方差矩阵或波束赋形向量，使通信用户的 SINR、速率或误码率约束与感知方向的波束图、照射功率、角度可分辨性等需求同时得到满足。此类问题通常可建模为多目标优化、约束优化或鲁棒优化问题，并常采用半正定松弛（SDR）、交替优化（AO）、加权和法、罚函数法以及分块坐标下降等方法求解。MIMO 联合设计的重要意义在于，它把通信与感知的冲突和协同关系明确地转化为统一优化框架，为后续大规模 MIMO、毫米波波束管理和网络级协同感知奠定了方法基础。
    \item \textbf{面向链路级折中的联合性能优化：}除了具体的波形和阵列设计之外，第二阶段还逐步形成了“通信--感知性能折中”这一核心研究主题\cite{Liu2020JointRadarComm,Liu2022ISAC}. 典型工作包括建立速率与检测性能之间的折中关系、刻画通信误码率与雷达参数估计精度之间的约束关系，以及利用帕累托边界、多目标优化和约束最优化方法分析双功能系统的可达性能区域。正是在这一类研究中，ISAC 开始形成自身独特的性能分析语言和方法论框架。
\end{itemize}
总体而言，第二阶段的意义在于：它首次在物理层和链路层层面系统性地证明了通信与感知并非只能“共存”，而是可以在同一信号、同一平台和同一优化框架下实现真正的双功能融合\cite{Liu2020JointRadarComm,Liu2022ISAC,Hassanien2015dual}. 通过这一阶段的发展，研究者逐步建立起一套围绕双功能波形、联合波束赋形、资源联合优化和通信--感知性能折中的基本理论、数学模型与算法工具，并验证了 OFDM、MIMO 等主流通信技术作为 ISAC 载体的可行性与有效性。然而，这一阶段的研究重心仍然主要停留在链路级和平台级，其对网络级协同、服务化感知和原生架构支持的关注仍然有限。也正因如此，随着 6G 愿景和网络智能化需求的提出，ISAC 技术进一步迈向了第三阶段，即面向网络与系统原生能力的内生通感一体化阶段。

\subsubsection{Stage 3: 6G Native Integrated Sensing and Communication Phase}
\label{subsubsec:ch1-stage3}

第三阶段代表了 ISAC 技术发展的前沿方向，也被广泛视为面向 6G 无线网络的重要愿景之一\cite{Liu2022ISAC,Gonzalez2024integrated,Saad2020Vision6G}。这一阶段最根本的特征在于：感知能力不再被看作附加在通信系统之外的外挂功能，而是被提升为无线网络的内生基本能力（native capability）。所谓“内生”，强调的是感知功能从网络体系结构设计之初就被纳入整体规划之中，并在标准协议、网络架构、资源调度、服务编排以及端到端业务支撑等多个层面获得原生支持。与前两个阶段相比，这一阶段的 ISAC 不再主要停留在物理层波形设计与链路级性能优化层面，而是进一步扩展为覆盖物理层、链路层、网络层乃至应用与服务层的系统级设计问题\cite{Liu2022ISAC,Zhang2020perceptive}。
从技术内涵上看，第三阶段的演进本质上体现为 ISAC 研究对象从“单链路双功能平台”向“网络化、服务化、智能化的感知通信系统”扩展。换言之，第二阶段解决的重点是“如何让一个平台同时通信和感知”，而第三阶段进一步关心的是“如何让一个网络原生地具备感知能力，并能够以服务化方式支撑复杂应用”\cite{Zhang2020perceptive,Cui2024integrated}. 在这一背景下，通信与感知的融合边界被显著拓展，网络级协同、计算协同、业务闭环与智能控制开始成为 ISAC 的核心组成部分。
围绕第三阶段，具有代表性的技术问题主要包括以下几个方面：
\begin{itemize}
    \item \textbf{网络级协同感知（Network-Level Cooperative Sensing）：}如何协调蜂窝网络中多个基站、接入点或感知节点的观测资源，实现更大范围、更高鲁棒性和更高空间分辨率的协同感知\cite{Zhang2020perceptive,Cui2024integrated}？在这一过程中，多节点之间的时间同步、频率同步、相位一致性、观测数据回传、分布式融合处理和联合目标跟踪等问题都将显著影响系统性能。
    \item \textbf{蜂窝 ISAC 组网（Cellular ISAC Networking）：}如何在蜂窝网络拓扑下部署和管理感知功能，使感知覆盖、感知分辨能力与通信覆盖、用户容量和干扰管理之间形成协调\cite{Liu2022ISAC,Zhang2020perceptive}？这一问题涉及感知小区与通信小区的关系、感知任务触发机制、空口资源配置方式以及接入网侧的功能拆分等多个方面。
    \item \textbf{端--边--云协同处理（Device--Edge--Cloud Collaborative Processing）：}由于感知原始数据通常规模较大、时效性要求较高，如何在终端、边缘节点和云平台之间合理分配感知数据的处理、压缩、融合与决策任务，成为第三阶段的重要系统问题\cite{Gonzalez2024integrated,Saad2020Vision6G}。这一问题不仅关系到感知精度，也直接影响网络回传负载、系统时延和能耗表现。
    \item \textbf{业务闭环支撑（Service-Loop Support）：}如何将感知结果实时反馈给通信、计算与控制系统，形成“感知--传输--计算--决策--控制”的闭环业务链路，是第三阶段区别于前两个阶段的重要特征之一\cite{Liu2022ISAC,Saad2020Vision6G}。在自动驾驶、工业控制、无人系统协同和数字孪生等场景中，感知不再只是一个独立输出，而是整个任务闭环中的起始环节和关键输入。
\end{itemize}

与前两个阶段相比，第三阶段的技术边界已明显从链路级融合扩展到系统级、网络级和业务级融合。具体而言，通感一体化的设计对象不再局限于单个基站、单个平台或单条链路，而是扩展到整个网络系统的统一优化。这种扩展至少体现在三个层面\cite{Zhang2020perceptive,Cui2024integrated}：第一，在资源管理层面，需要面向全网进行感知资源调度、多小区干扰与协同管理、感知覆盖规划以及跨节点观测协同；第二，在架构与协议层面，感知功能需要被纳入网络功能实体和控制面设计，例如在核心网或边缘侧引入感知任务管理、感知数据融合分析与服务编排相关的功能模块，并在接入网空口与控制协议中原生支持感知信号的发射、观测与反馈；第三，在业务层面，通信服务与感知服务不再沿两条彼此独立的管线运行，而是面向上层任务进行统一编排与交付，使应用能够按需请求“连接 + 感知 + 计算”的组合能力。
在这一阶段，ISAC 系统开始展现出前两个阶段难以实现的新型能力与系统价值。
\begin{itemize}
    \item \textbf{感知辅助通信（Sensing-Assisted Communication）：}利用感知结果反向提升通信系统性能。例如，通过获取用户位置、运动轨迹、环境遮挡状态和空间传播结构等信息，网络可以提前进行波束预测、链路切换准备、阻塞预警和资源调度优化，从而降低波束训练开销、提升链路稳定性并减少切换时延\cite{Liu2022ISAC,Zhang2020perceptive}。这一能力在毫米波和太赫兹频段中尤为重要，因为高频窄波束通信对空间状态信息高度敏感。
    \item \textbf{通信赋能感知（Communication-Empowered Sensing）：}利用通信网络已有的基础设施、信号资源和时空覆盖能力实现大规模、持续性环境感知\cite{Liu2022ISAC,Gonzalez2024integrated}。例如，蜂窝网络的多基站部署为分布式协同感知提供了天然的观测节点网络，持续存在的下行与上行通信信号可作为机会照射源，网络同步与回传机制则为多节点数据融合提供了基础条件。与传统独立感知系统相比，这种方式更有利于实现广域覆盖和基础设施复用。
    \item \textbf{感知驱动的网络智能（Sensing-Driven Network Intelligence）：}将感知获取的物理环境信息纳入网络智能决策过程，使资源调度、移动性管理、负载均衡、波束控制和服务质量保障不再仅依赖通信侧反馈，而是能够结合对物理世界状态的实时理解进行优化\cite{Saad2020Vision6G,Cui2024integrated}。这也使得 AI/ML 技术与 ISAC 之间形成更加紧密的耦合关系，从而推动网络从“自适应连接系统”进一步演化为“具备环境认知能力的智能系统”。
    \item \textbf{服务闭环（Service Loop）：}形成“感知--通信--计算--决策--控制--再感知”的完整闭环能力\cite{Saad2020Vision6G,Liu2022ISAC}。例如，在工业自动化场景中，网络侧感知系统发现设备运行异常后，可通过通信链路将观测信息快速传送至边缘计算节点，由后者完成状态评估和控制决策，再通过控制链路向执行设备下发调整指令，最后由感知功能继续监测控制效果并形成反馈。这种闭环能力意味着 ISAC 不再只是通信系统的功能扩展，而是成为支撑数字化、智能化和自动化应用的重要基础能力。
\end{itemize}
从技术演进的意义上看，第三阶段标志着 ISAC 从“链路级双功能技术模块”迈向“网络级原生能力框架”的根本性转变\cite{Liu2022ISAC,Zhang2020perceptive,Cui2024integrated}。在这一阶段，ISAC 不再只是物理层研究中的信号处理问题，而是同时成为网络架构设计、协议演进、边缘智能、标准化推进和产业生态构建所共同关注的方向。与此同时，这一阶段也提出了大量尚待深入研究的开放问题，例如网络级 ISAC 的基础性能极限、多节点协同感知的最优任务分配机制、感知数据的安全与隐私保护、感知即服务（Sensing as a Service）的业务模式与资源计费方法等\cite{Gonzalez2024integrated,Cui2024integrated}。
从标准化与产业演进角度看，ISAC 也正在从学术研究逐步走向工程体系建设。近年来，3GPP 已围绕 ISAC 相关能力开展研究讨论，并在 Rel-19 阶段启动相关研究工作；同时，面向 IMT-2030 的 6G 愿景研究中，也已将感知与通信融合视为未来无线网络的重要发展方向之一\cite{IMT2030Framework,ThreeGPPRel19ISAC}. 这些进展表明，第三阶段并非抽象愿景，而是正在逐步进入可研究、可实现、可标准化推进的工程实践阶段。

\subsection{The core driving force behind the three-stage evolution}
\label{subsec:ch1-evolution-driving-force}

ISAC 技术之所以沿着前述三阶段路径不断演进，并不是若干孤立技术方案偶然叠加的结果，而是由一系列长期存在且不断增强的结构性驱动力共同推动的。这些驱动力既包括频谱与硬件资源层面的客观约束，也包括未来网络能力定位、应用场景需求和系统智能化趋势所带来的功能牵引。理解这些驱动力，有助于从更高层面把握 ISAC 技术从“共存”走向“融合”、再从“链路级融合”走向“网络级内生能力”的内在逻辑。

\subsubsection{Spectral efficiency and resource reuse requirements}
\label{subsubsec:ch1-driving-force-spectrum}

频谱资源的稀缺性是推动 ISAC 发展的最直接驱动力之一。随着全球移动数据业务持续增长，无线网络对频谱资源的需求不断攀升，而适合移动通信使用的优质连续频谱，尤其是 Sub-6~GHz 频段，已日趋紧张\cite{ITU_R_RR2024,miit_radio_freq_2026_en,Saad2020Vision6G}。与此同时，为满足更高吞吐量和更大系统容量需求，通信系统不断向毫米波乃至太赫兹频段拓展，而这些频段与雷达、遥感和无线电定位等感知系统的工作频率存在越来越明显的重叠或邻频关系\cite{Liu2022ISAC}。正是在这一背景下，通信与感知之间的频谱冲突首先推动了第一阶段中频谱共存与干扰管理技术的发展。

然而，共存只能在一定程度上缓解冲突，却难以从根本上提升频谱使用效率。其根本原因在于，在共存框架下，通信系统和感知系统虽然共享频谱，但本质上仍然是竞争关系：一个系统性能的提升往往伴随着另一个系统可用资源的压缩或干扰水平的上升\cite{Liu2020JointRadarComm,Chiriyath2017RadarCommCoexistence}。正因如此，单纯的共存机制无法充分释放共享频谱的潜力。相比之下，如果通信与感知能够共享同一发射信号、同一时频资源乃至同一硬件平台，则同一份物理资源便可以同时服务于信息传输与环境感知两类功能，从而带来超越简单分时/分频共享的资源复用收益。这种从“频谱共享”走向“信号与功能共享”的演进，构成了第二阶段联合信号设计的重要根本动因\cite{Liu2022ISAC,Hassanien2015dual}。

更深层次地看，随着 6G 网络需要同时支撑通信、感知、定位、计算与控制等多种能力，传统“一种功能占用一套专属资源”的系统设计范式正在变得不可持续\cite{Gonzalez2024integrated,Saad2020Vision6G}。未来网络中的时间、频率、空间和功率等资源必须服务于多重功能目标，这意味着系统设计理念必须从“功能专属”逐步转向“资源复用、功能协同和统一调度”。ISAC 正是这一演进趋势在无线空口与网络体系中的典型体现。

\subsubsection{Hardware integration and reconfigurable platform development}
\label{subsubsec:ch1-driving-force-hardware}

如果说频谱压力提供了 ISAC 发展的外部牵引，那么硬件平台能力的提升则为其实现提供了关键物理基础。通信与感知之所以能够从“理论上可结合”逐步走向“工程上可实现”，一个重要原因就在于现代无线硬件平台正在变得更加可重构、宽带化、阵列化和软件可编程。

首先，软件定义无线电（Software Defined Radio, SDR）以及通用化基带平台的发展，使得同一套射频与基带硬件能够通过软件配置支持不同的信号体制和处理流程，从而显著降低了通信与感知功能在同一平台上实现的门槛\cite{Mitola1995SDR}. 这种可编程性意味着系统不再需要为不同功能部署完全独立的硬件链路，而可以通过软件升级、参数重构和功能切换实现多任务复用，这为第二阶段的双功能平台设计提供了基础条件。
其次，大规模 MIMO 和相控阵技术的发展，为 ISAC 提供了丰富的空间自由度。通过配置大量天线单元，系统能够在空间域同时形成多个具有不同方向、宽度和功率分布特征的波束，从而为通信用户提供高质量定向覆盖，也为感知目标提供可控照射与角度解析能力\cite{Liu2020JointRadarComm,Li2023integrated}。正是这种空间可塑性，使得联合波束赋形和空间资源复用成为可能，也使第二阶段中大量基于 MIMO 的联合设计方案具有了现实可行性。
再次，毫米波和太赫兹通信技术天然具备与感知需求高度契合的物理特性。高频段系统通常拥有更大的可用带宽，因此能够提供更高的距离分辨率。以经典单基地体制为例，其距离分辨率近似满足:
\[
\Delta R = \frac{c}{2B},
\]
其中 $c$ 为光速( m/s )，$B$ 为有效信号带宽( Hz )\cite{SkolnikRadarBook}. 这意味着带宽越大，目标之间在距离维上的可分辨能力越强。与此同时，毫米波和太赫兹系统通常依赖大规模阵列与窄波束传输，这又天然有利于提高角度分辨率和空间聚焦能力\cite{Rappaport2019WirelessTHz}. 因此，从物理层属性来看，高频通信系统本身就具备承担高精度感知任务的潜力，这也是 ISAC 在毫米波和太赫兹频段特别受到关注的重要原因。

还需要指出的是，在较低频段中，传统通信设备与传统雷达设备在射频架构、带宽需求和阵列形态上存在较大差异，因此硬件共享往往较为困难；而在毫米波及更高频段，通信与感知系统在前端实现上却呈现出越来越明显的趋同性，例如都倾向于采用宽带射频链路、相控阵结构、高速模数/数模转换器以及紧耦合的数字波束成形架构\cite{Liu2022ISAC,Rappaport2019WirelessTHz}。这种硬件架构的趋同，显著降低了通信与感知共平台实现的工程障碍，推动了 ISAC 从“概念可行”走向“系统可实现”。
此外，新型电磁可重构器件的发展也为 ISAC 打开了新的设计空间。以可重构智能表面（Reconfigurable Intelligent Surface, RIS）为代表的可编程电磁环境技术，可以通过调节反射单元的幅相响应，主动塑造无线传播环境，从而同时服务于通信链路增强与感知覆盖扩展\cite{Chepuri2023integrated}. 在 ISAC 场景下，RIS 不仅能够改善覆盖盲区、增强感知回波或构造有利的传播路径，还能够为通信与感知的联合优化提供额外自由度，因此成为近年来的重要研究热点之一。

\subsubsection{Network intelligence and service loop requirements}
\label{subsubsec:ch1-driving-force-network}

除了资源与硬件层面的推动，移动通信网络自身定位的变化同样是 ISAC 演进的重要驱动力。5G 及以前的无线网络，其核心价值主要在于提供高质量连接；而面向 6G 的愿景中，网络将进一步演进为集连接、感知、计算、控制与智能于一体的综合服务平台\cite{Gonzalez2024integrated,Saad2020Vision6G}。在这一转变过程中，越来越多新兴应用对“感知--通信--计算--决策--控制”闭环服务的需求不断增强，从而使感知能力从外围辅助模块逐步上升为网络体系中的关键内生能力。
例如，在自动驾驶场景中，系统需要在毫秒级时间尺度上完成环境检测、状态共享、路径规划和控制执行；在工业互联网与智能制造场景中，系统需要完成“监测--传输--分析--调整”的实时闭环；在远程控制与数字化运维场景中，网络还需要持续感知物理设备和环境状态，以支撑高可靠反馈控制\cite{Liu2022ISAC,Saad2020Vision6G}. 在这些闭环中，感知并不是一个可有可无的附加输出，而是整个业务流程的起始环节和关键输入。如果感知功能完全依赖于网络之外的独立系统，则感知数据回传时延、接口复杂度、时钟同步和异构数据融合等问题都可能成为制约闭环性能的瓶颈。因此，将感知能力内生到网络体系之中，成为实现高效闭环服务的重要前提。

与此同时，随着 AI/ML 技术在无线网络优化中的广泛应用，网络决策越来越依赖于对环境状态的准确理解。例如，智能波束管理需要获取用户位置与运动方向，智能资源调度需要把握业务流量在时间与空间上的分布规律，智能干扰管理需要识别潜在干扰源的位置与活动状态\cite{Saad2020Vision6G}. 传统网络主要依靠信道反馈、测量报告和历史统计信息进行优化，这些信息虽然重要，但往往在环境可解释性、实时性和前瞻性方面存在局限。而 ISAC 通过在网络中原生引入环境感知能力，使网络能够获得更丰富、更直接的物理世界信息，从而为网络智能化提供高价值输入。这一趋势直接推动了第三阶段中“感知驱动网络智能”的形成。

\subsubsection{6G Vision Requirements for Intrinsic Capabilities}
\label{subsubsec:ch1-driving-force-6g-vision}

从更宏观的角度看，6G 发展愿景本身就对无线网络提出了超越传统连接能力的更高要求。研究界普遍认为，面向未来的网络不应仅仅提供高速传输和广域连接，还应具备对环境、目标、用户和业务状态的实时感知、场景理解与协同控制能力\cite{Saad2020Vision6G,IMT2030Framework}. 在这一过程中，超高可靠低时延、原生智能、数字孪生、低空经济以及空天地海一体化等新型场景的共同需求，正在推动通信系统由“单一连接能力”向“通信、感知、计算与智能深度融合”的内生能力体系演进，而 ISAC 正是支撑这一演进的重要技术基础\cite{Gonzalez2024integrated,Liu2022ISAC}。

首先，超高可靠与低时延业务对网络提出了更强的预测性与确定性要求。传统无线网络往往依赖链路质量下降后的被动调整机制，难以满足未来业务对毫秒级甚至更低时延以及接近零中断服务的要求。ISAC 通过在通信过程中同步获取用户位置、移动轨迹、遮挡变化、环境反射体分布以及潜在干扰信息，使网络能够更早感知信道和拓扑变化，并提前执行波束管理、链路切换与资源配置\cite{Liu2022ISAC,Zhang2020perceptive}. 这意味着网络不再只是“事后补救”，而开始具备“提前感知、提前决策、提前优化”的能力。

其次，原生智能的提出进一步强化了网络对环境信息的依赖。智能网络不仅需要传输数据，还需要基于环境状态和业务状态进行自主学习、推理和优化。而其中最关键的一类输入，正是来自真实物理世界的动态感知信息\cite{Saad2020Vision6G}. 通过将感知能力嵌入网络基础设施，ISAC 能够持续获取用户行为、场景变化、设备状态和空间分布等多维信息，为网络 AI 模型提供训练样本与实时推理输入，从而支撑资源调度、业务预测、异常检测、移动性管理和多任务协同等智能功能。

再次，数字孪生愿景进一步放大了网络对实时感知和持续更新能力的需求。数字孪生的本质是对物理世界进行高精度、动态化和可交互的数字映射，而这种映射不可能仅依赖离线建模或静态数据库，而必须依赖持续、可靠、广覆盖的感知机制不断进行刷新与校准\cite{IMT2030Framework,Khan2022digital}. ISAC 能够利用现有通信网络基础设施，在提供连接服务的同时采集目标位置、运动状态、环境结构和场景变化信息，从而为数字孪生提供实时数据入口。也正因如此，ISAC 可以被看作未来网络支撑数字世界与物理世界耦合的重要感知基础。

此外，低空经济以及空天地海一体化等新型场景，也对网络提出了更强的“边连接、边感知、边决策”能力需求\cite{Saad2020Vision6G,IMT2030Framework}. 这些场景通常具有高度动态、广域分布、节点异构、传播环境复杂和链路状态快速变化等特点，仅依赖传统通信机制很难实现高质量、稳定且安全的服务保障。例如，在低空飞行器网络中，系统不仅要完成控制和业务数据传输，还需要实时掌握飞行器位置、速度、航迹演化及周边空域态势；在空天地海一体化网络中，卫星、无人机、地面基站和海上平台之间的协同，也需要对节点状态、环境条件和目标活动进行联合感知与动态调度。在这些场景下，ISAC 使网络具备环境理解和风险预警能力，从而增强复杂场景下的自治组网、动态覆盖与协同控制能力。

总体而言，6G 愿景对网络提出的并不是某一单一性能指标的局部提升，而是要求网络具备一种面向物理世界的综合内生能力，即能够实时感知环境、理解场景、预测变化并自主优化\cite{Liu2022ISAC,Gonzalez2024integrated,Saad2020Vision6G}. 这些要求共同指向同一个发展趋势：未来无线网络必须将感知能力深度嵌入系统架构之中，使其成为与连接能力同等重要的基础能力。也正是在这一意义上，ISAC 不仅是 6G 中的一个关键技术方向，更是支撑 6G 内生能力形成和网络形态跃迁的重要基础。

\subsection{Summary of Technical Boundaries and Core Characteristics at Each Stage}
\label{subsec:ch1-evolution-summary}

在前文分别讨论了 ISAC 三阶段演进逻辑及其背后的主要驱动力之后，本节进一步从技术目标、设计层级和评价指标三个维度，对频谱共存阶段、联合设计阶段以及内生一体化阶段进行系统性总结，以帮助读者从整体上把握不同阶段之间的联系、差异及其递进关系。需要强调的是，这种分阶段总结并不是为了将不同研究方向割裂开来，而是为了提炼出通信与感知融合深度不断增强的技术主线。

\subsubsection{Technical target differences}
\label{subsubsec:ch1-target-differences}

从技术目标上看，三个阶段体现出从“减少冲突”到“共享资源”，再到“形成系统级闭环能力”的递进逻辑。
\begin{itemize}
    \item \textbf{第一阶段：减少冲突与维持兼容运行。}  
    在频谱共存阶段，通信系统与雷达/感知系统通常仍然是两个彼此独立的系统。此时的核心目标是在共享频谱或邻频工作条件下尽量降低两类系统之间的相互干扰，使双方性能退化控制在可接受范围内。换言之，这一阶段追求的并不是功能层面的协同增益，而是兼容性与可共存性，其设计哲学可以概括为“和平共处”。
    \item \textbf{第二阶段：共享资源并兼顾双功能性能。}  
    在联合设计阶段，研究重点转向“如何让同一平台、同一信号或同一组资源同时服务于通信与感知”\cite{Zhang2021ISAC6G,Hassanien2015dual}。因此，该阶段的核心目标不再只是降低冲突，而是通过双功能波形设计、联合波束赋形和资源共享机制，在给定约束条件下兼顾通信与感知两类性能需求，并寻求二者之间的最优折中。这一阶段的设计哲学更接近“合作共赢”。
    \item \textbf{第三阶段：构建网络级原生能力与业务闭环。}  
    在内生一体化阶段，感知能力被纳入 6G 及未来无线网络的原生能力体系之中，并与通信、计算、控制和智能功能共同支撑端到端业务流程\cite{Saad2020Vision6G,Zhang2020perceptive,Cui2024integrated}。因此，这一阶段的核心目标不再局限于链路级双功能实现，而是进一步上升为形成面向应用任务的系统闭环能力，即通过“感知--通信--计算--决策--控制”的联合支撑，为自动驾驶、工业控制、数字孪生和低空协同等场景提供整体价值。其设计哲学可概括为“浑然一体”。
\end{itemize}

\subsubsection{Design level differences}
\label{subsubsec:ch1-design-level-differences}

三个阶段在设计所关注的层级上，也呈现出由底层向高层、由局部向全局逐步扩展的趋势。
\begin{itemize}
    \item \textbf{第一阶段：频谱与干扰管理层。}  
    在第一阶段，设计关注点主要集中在频谱使用与干扰控制层面。其核心问题是如何通过频率规划、保护带设计、功率控制、方向隔离、认知频谱接入以及干扰抑制等手段，实现通信系统与雷达系统的兼容运行\cite{Chiriyath2017RadarCommCoexistence,Blunt2018RadarSpectrumSharing}。这一阶段的设计对象主要是“资源冲突”本身，而非统一系统架构。
    \item \textbf{第二阶段：波形、波束与资源联合优化层。}  
    在第二阶段，设计关注点进一步扩展到物理层和链路层的联合设计问题，包括双功能波形构造、MIMO 联合波束赋形、发射协方差矩阵优化以及时--频--空--功率资源的统一调度\cite{Liu2020JointRadarComm,Zhang2021ISAC6G,Sturm2011WaveformDesign}。这一阶段的设计对象已从“避免冲突”转变为“形成统一信号与统一资源优化框架”，技术手段主要包括双功能波形设计、通信--感知性能折中分析、多目标优化和约束优化等。
    \item \textbf{第三阶段：网络、架构与业务层。}  
    在第三阶段，设计关注点进一步提升到网络体系结构、控制面协议、服务编排与业务闭环层面\cite{Gonzalez2024integrated,Zhang2020perceptive,Cui2024integrated}。其核心问题不再只是某一个节点如何同时通信与感知，而是整个网络如何原生支持感知功能，包括多节点协同感知的任务调度、感知数据的回传与融合、网络功能实体中的感知管理模块、端--边--云协同处理以及面向应用的服务接口设计等。此时，ISAC 已从一种链路级技术问题逐步上升为一种系统级和业务级能力构建问题。
\end{itemize}

\subsubsection{Evaluation metric differences}
\label{subsubsec:ch1-metric-differences}

随着设计目标和系统层级的变化，三个阶段在性能评价上关注的核心指标也表现出显著差异。
\begin{itemize}
    \item \textbf{第一阶段：以互扰抑制与兼容性评估为主。}  
    在频谱共存阶段，评价重点主要围绕“共存条件下各系统性能退化了多少”这一问题展开。典型指标包括干扰噪声比（INR）、干扰条件下通信链路的信干噪比（SINR）、误码率与吞吐量，以及雷达/感知系统在干扰存在下的检测概率、虚警概率、估计误差和干扰保护裕量等\cite{Chiriyath2017RadarCommCoexistence,Blunt2018RadarSpectrumSharing}。这些指标反映的是系统在互扰环境中的鲁棒性和可共存能力。
    \item \textbf{第二阶段：以通信--感知性能折中为主。}  
    在联合设计阶段，评价重点转向“在给定资源和约束下，通信性能与感知性能能够达到怎样的联合最优折中”\cite{Liu2020JointRadarComm,Liu2022ISAC}。因此，典型指标除了通信侧的速率、SINR、误码率外，还包括感知侧的检测概率、估计精度、CRLB、距离/速度/角度分辨率以及波束图增益等；进一步地，还会通过帕累托边界、速率--估计精度区域（例如 rate--CRLB region）或多目标最优前沿来刻画双功能系统的性能极限\cite{Zhang2021ISAC6G}。这一阶段评价的核心问题是“如何量化协同收益与性能折中”。
    \item \textbf{第三阶段：以系统级综合收益与业务闭环效果为主。}  
    在内生一体化阶段，单纯的链路级指标已不足以完整评价系统价值，性能评估需要进一步扩展到网络级和业务级层面\cite{Zhang2020perceptive,Cui2024integrated}。除了通信和感知的基础链路指标之外，还需要考虑网络级感知覆盖率、感知容量、多节点协同感知收益、端到端业务时延、服务可用性、计算与回传开销、资源利用效率以及闭环任务完成质量等更综合的指标。此时，评价的核心问题已从“某个信号或某条链路表现如何”转变为“ISAC 作为网络原生能力，为系统运行和应用任务带来了多大综合价值”。
\end{itemize}
为了更直观地展示三阶段之间的差异，表~\ref{tab:stage-comparison} 对其主要技术特征作进一步总结。

\begin{table}[!ht]
\centering
\small
\renewcommand{\arraystretch}{1.2}
\setlength{\tabcolsep}{4pt}
\begin{tabularx}{\textwidth}{|>{\centering\arraybackslash}p{2.5cm}|
                             >{\centering\arraybackslash}X|
                             >{\centering\arraybackslash}X|
                             >{\centering\arraybackslash}X|}
\hline
\textbf{Comparison Dimension} & 
\textbf{Stage 1: Spectrum Coexistence} & 
\textbf{Stage 2: Joint Design} & 
\textbf{Stage 3: Native Integration} \\
\hline
\textbf{Core objective} & 
Conflict mitigation & 
Resource sharing with dual functionality & 
System-level closed-loop capability \\
\hline
\textbf{Design philosophy} & 
Peaceful coexistence & 
Cooperative dual-function design & 
Native and holistic integration \\
\hline
\textbf{Design level} & 
Spectrum/interference layer & 
Waveform/beam/resource layer & 
Network/\allowbreak architecture/\allowbreak service layer \\
\hline
\textbf{System relation} & 
Coordination of two independent systems & 
Dual-functional design on one platform & 
Native capability of one network \\
\hline
\textbf{Core metrics} & 
Interference suppression and compatibility & 
Communication--sensing tradeoff & 
System-wide utility and service value \\
\hline
\textbf{Representative techniques} & 
Spectrum sharing, interference management & 
DFRC waveform, joint beamforming & 
Cellular sensing, cooperative networking, service loop \\
\hline
\textbf{Maturity level} & 
Relatively mature & 
Rapidly developing & 
Early exploration stage \\
\hline
\end{tabularx}
\caption{Systematic comparison of the technical characteristics of the three stages}
\label{tab:stage-comparison}
\end{table}

最后需要指出的是，上述三阶段并不应被机械地理解为严格的时间先后序列。事实上，在当前研究与工程实践中，三个阶段的问题往往是并行推进的：一方面，频谱共存问题在现实部署中依然长期存在；另一方面，链路级联合设计仍然是当前 ISAC 研究的核心内容之一，而网络级原生一体化也正在逐步展开\cite{Liu2022ISAC,Gonzalez2024integrated,Cui2024integrated}。因此，三阶段划分更准确地说反映的是技术融合深度、系统集成层级和能力边界的递进关系，而不是一个严格线性的历史时间表。许多现实系统往往同时具有不同阶段的技术要素，这也是理解 ISAC 复杂性和开放性的一个重要方面。

\subsection{Reflective Questions}
\label{subsec:ch1.1-reflect-questions}

本节思考题分为“概念理解”“对比分析”“深度思考与论述”和“开放探索”四类。建议读者在完成本章学习后独立作答，其中部分题目也可作为课堂讨论、小组研讨或课程论文选题。

\paragraph{一、概念理解类}
\begin{enumerate}
    \item 请用自己的语言解释：为什么说 ISAC 不能简单理解为“通信 + 雷达”的物理拼接？试从信号层、资源层、硬件层、协议层、网络层和业务层等维度，举例说明“深层耦合”与“简单拼接”的区别。

    \item 在通信系统中，无线信道通常被视为需要“克服”的不利因素；而在雷达/感知系统中，信道（包含目标响应）却是需要“分析”的信息载体。请结合这一本质差异，讨论 ISAC 系统在信道建模时面临的独特挑战：如何使同一信道模型同时服务于通信功能和感知功能？

    \item 本章提到，ISAC 中通信与感知之间存在“双向赋能”关系。请分别解释“感知辅助通信”和“通信赋能感知”的含义，并各举一个具体应用场景进行说明。若去掉双向赋能中的任一方向，系统性能可能会受到怎样的影响？

    \item 请辨析以下术语之间的联系与区别：ISAC、JCAS、JRC、DFRC、RadCom。为什么本教材选择以 ISAC 作为统一术语？
\end{enumerate}

\paragraph{二、对比分析类}
\begin{enumerate}
    \setcounter{enumi}{4}
    \item 请从发射信号特性、接收处理方式、目标对象、性能评价指标以及信道角色定位等维度，对比通信系统与雷达/感知系统的核心差异，并总结这些差异给 ISAC 统一设计带来的主要矛盾。

    \item 对比 ISAC 演进的三个阶段——频谱共存阶段、联合设计阶段和 6G 内生一体化阶段——在以下方面的差异：
    \begin{enumerate}
        \item 核心技术目标；
        \item 设计关注的层级；
        \item 通信与感知的系统关系；
        \item 典型性能评价指标。
    \end{enumerate}
    并进一步讨论：这三个阶段之间是否存在严格的时间先后关系？为什么？

    \item 在频谱共存阶段，典型实现方式包括时分共享、频分隔离、共存感知和机会式频谱复用。请分析这四种方式各自的优缺点，并解释为什么这些方式在频谱效率上的提升空间有限，从而推动技术向第二阶段演进。
\end{enumerate}

\paragraph{三、深度思考与论述类}
\begin{enumerate}
    \setcounter{enumi}{7}
    \item 本章指出，5G 及以前的移动通信网络的核心角色主要是“信息传输管道”，而 6G 网络正在向“连接、感知、计算、控制”一体化的综合服务平台演进。请结合一个你熟悉的垂直行业应用场景（如自动驾驶、智能工厂、低空无人机管理等），详细回答：
    \begin{enumerate}
        \item 该场景中的通信需求与感知需求分别是什么？
        \item 为什么传统的“外挂式”独立感知方案难以满足这些需求？
        \item ISAC 如何在该场景中实现“感知 $\rightarrow$ 通信 $\rightarrow$ 计算 $\rightarrow$ 决策 $\rightarrow$ 控制”的业务闭环？
    \end{enumerate}

    \item 毫米波/太赫兹频段的通信系统天然具备大带宽和窄波束特性。请从距离分辨率
    \[
        \Delta R = \frac{c}{2B}
    \]
    和角度分辨率两个角度，解释为什么高频段通信系统天然适合承担感知功能。同时讨论高频段 ISAC 系统面临的独特技术挑战，如路径损耗、波束对准、硬件复杂度与能耗等。

    \item 通信信号（如 OFDM 信号）通常具有随机性，因为其承载了未知的数据符号；而传统雷达信号通常是确定性已知的，以便接收端进行匹配滤波。在 ISAC 系统中，若采用 OFDM 通信信号同时执行感知任务，而感知接收端无法完整获得瞬时发射数据符号，接收端应如何处理“参考信号未知或部分未知”的问题？这一问题对 ISAC 的系统架构设计有何影响？

    \item 本章将 ISAC 的演进驱动力概括为四个方面：频谱效率需求、硬件平台发展、网络智能化趋势和 6G 愿景。请分析：
    \begin{enumerate}
        \item 这四个驱动力之间是否存在内在关联？如存在，请用图示或文字说明它们之间的逻辑关系；
        \item 你认为其中哪个驱动力最为根本？为什么？
        \item 除了这四个驱动力之外，你认为还有哪些因素可能推动 ISAC 技术的演进？
    \end{enumerate}

    \item 在 ISAC 第二阶段的联合波束成形设计中，系统通常需要同时满足通信用户的 SINR 需求和感知目标方向的照射功率需求。请从直观上分析：当通信用户和感知目标在空间上位于同一方向时，联合波束成形设计会更容易还是更困难？当它们在空间上近似正交时又如何？请给出你的分析与理由。
\end{enumerate}

\paragraph{四、开放探索类}
\begin{enumerate}
    \setcounter{enumi}{12}
    \item 3GPP 和 ITU-R 已将感知能力纳入 6G 相关研究议程。请查阅最新的标准化文献或白皮书（作答时注明版本或发布时间），回答：
    \begin{enumerate}
        \item 当前 ISAC 标准化研究的范围和重点方向是什么？
        \item ISAC 在标准化过程中面临的主要争议或难点有哪些？
        \item 你认为 ISAC 从标准化走向大规模商用部署，还需要克服哪些关键障碍？
    \end{enumerate}

    \item 本章提到，ISAC 的感知对象非常广泛，不仅包括传统的运动目标检测，还包括生命体征监测、手势识别等。请选择一种你感兴趣的非传统感知应用，讨论：
    \begin{enumerate}
        \item 该应用对感知精度、时延和覆盖范围有哪些具体需求？
        \item 利用通信信号（如 Wi-Fi、5G NR）实现此类感知的基本原理是什么？
        \item 这类感知应用可能引发哪些隐私与安全方面的担忧？
    \end{enumerate}

    \item 设想你是一家移动运营商的 CTO，正在评估是否在现有 5G 基站上引入 ISAC 能力。请从以下角度进行分析：
    \begin{enumerate}
        \item 引入 ISAC 功能可能带来哪些新的业务收入来源？
        \item 需要对现有基站的硬件与软件进行哪些升级？
        \item 部署 ISAC 功能面临哪些技术与商业风险？
        \item 你会建议优先在哪些场景中部署 ISAC？为什么？
    \end{enumerate}
\end{enumerate}

\noindent\textbf{提示：} 以上思考题按照“概念理解—对比分析—深度论述—开放探索”四个层次进行设计。建议读者在完成本章学习后独立作答，其中部分题目适合作为课堂讨论、小组研讨或课程作业选题。
